\documentclass{article} % For LaTeX2e
\usepackage{iclr2027_conference,times}

\usepackage{hyperref}
\usepackage{url}
\usepackage{amsmath}
\usepackage{amssymb}
\usepackage{booktabs}
\usepackage{multirow}
\usepackage{array}
\usepackage[dvipsnames]{xcolor}
\usepackage{enumitem}

% ---------------------------------------------------------------------------
% Plain copy: no venue running head, no author block.  Comment out the \def
% below to restore the ICLR submission furniture ("Under review as a conference
% paper at ICLR 2027" and the anonymous-author line).
% ---------------------------------------------------------------------------
\def\PLAINCOPY{}
\ifdefined\PLAINCOPY
  \makeatletter
  \def\@maketitle{\vbox{\hsize\textwidth
    {\LARGE\sc \@title\par}
    \lhead{}%
    \vskip 0.3in minus 0.1in}}
  \makeatother
  \renewcommand{\headrulewidth}{0pt}
\fi

% ---------------------------------------------------------------------------
% Working-copy markup.  Everything typeset in red is an open item that is
% itemised in Appendix A (working notes).  Comment out the two lines below and
% the \workingnotes switch at the bottom of this file to produce a clean copy.
% ---------------------------------------------------------------------------
% Uncomment the next line to build the clean submission copy: the red markup and
% the working-notes appendix are dropped, and \chknote text disappears entirely.
%\def\CLEANCOPY{}

\newcommand{\red}[1]{\textcolor{red}{#1}}
\newcommand{\opentag}[1]{\textcolor{red}{\textsuperscript{[\ref{#1}]}}}
\newcommand{\chknote}[2]{\red{[\textsc{check} \ref{#2}: #1]}}
\newcommand{\pend}[2]{\red{#1}\opentag{#2}}

\ifdefined\CLEANCOPY
  \renewcommand{\red}[1]{#1}
  \renewcommand{\opentag}[1]{}
  \renewcommand{\chknote}[2]{}
  \renewcommand{\pend}[2]{#1}
\fi

\newcommand{\Lstruct}{\ensuremath{L_{\mathrm{struct}}}}
\newcommand{\Lseq}{\ensuremath{L_{\mathrm{seq}}}}
\newcommand{\dstruct}{\ensuremath{d_{\mathrm{struct}}}}
\newcommand{\angstrom}{\text{\AA}}

\title{When the Metric Goes Up:\\ A Training-Distribution Control for\\ Interpretability Evaluation}

% Author names are suppressed automatically in submission mode; the block below
% is used only once \iclrfinalcopy is uncommented.
\author{Wenjun Gao \\
Department of ... \\
University College London \\
\texttt{zcbtwga@ucl.ac.uk} \\
\And
Alex Fedorec \\
Department of ... \\
University College London \\
\texttt{email}
}

%\iclrfinalcopy % Uncomment for camera-ready version, but NOT for submission.

\begin{document}

\maketitle

\begin{abstract}
Protein language models (pLMs) are transformers over amino-acid sequences. Models trained on
sequence alone encode structure and function, and SAE-based interpretability metrics are used
to measure how well a model's features capture these properties. Such metrics are typically
validated by showing that the score is large and survives multiple-comparison correction. We
show that this is insufficient: a metric of this kind can rise when the relation it names is
removed from the training data. We define a co-activation statistic, \Lstruct, that scores
whether SAE features cluster on residues in spatial contact, and evaluate it on a controlled
pair of masked- and causal-objective models sharing initialisation, corpus and batch order.
The masked arm scores higher at every depth. We subject this result to six checks --- covering
dictionary training, reconstruction quality, amino-acid composition and architecture. It
survives all of them. The inference is nonetheless wrong. Retrained on a corpus with residue
order destroyed, \Lstruct{} rises in every model-by-depth combination tested. An untrained
network matches it. Direct probes for secondary structure and burial favour the opposite arm
in every comparison. A single amino-acid indicator outscores every learned feature in five of
six conditions. A published structural statistic, reimplemented from its Methods, fails the
same control. What survives is a training-invariant floor that neither corpus nor budget
moves.
\end{abstract}

\ifdefined\CLEANCOPY\else
\red{\emph{Working copy.} Red text marks an open item: a value to check against the code, or a
claim a queued run may change. Each carries a tag pointing to Appendix~\ref{app:working}, which
is not part of the paper. The appendices cited in the text (A--P) are not included here: that
material was lost with the drive and is being recovered. The two figures are described in place.}
\fi

\section{Introduction}

Protein language models (pLMs) trained on amino-acid sequences encode structure and function,
from residue contacts and secondary structure to binding sites and stability
\citep{rives2021biological,rao2021transformer}, though what underlies this is debated;
\citet{zhang2024protein} argue that pLMs capture coevolutionary statistics rather than protein
physics. Sparse autoencoders (SAEs) are one instrument for asking what a model encodes,
decomposing activations into a large dictionary of sparse features, each scored against protein
annotations or structures \citep{simon2025interplm,adams2025mechanistic,gujral2025sparse}.

Two kinds of score are in use, and they differ in what can validate them.
\emph{Annotation-referenced} scores compare a feature's activations against an externally
defined label, such as a Swiss-Prot site or a secondary-structure class, using F1 or AUROC
\citep{simon2025interplm,adams2025mechanistic}. They have a reference value: a perfect score is
defined, and the label exists independently of the feature. \emph{Relation-scored} statistics
ask instead whether the residues a feature fires on stand in a specified geometric relationship
to each other, such as spatial proximity within a distance cutoff, with the null generated by
permuting the feature's own activations within each protein \citep{simon2025interplm}. No
per-feature ground truth exists and no value is defined as correct, so the score is accepted on
indirect grounds: it is large and it survives multiple-comparison correction.

InterPLM \citep{simon2025interplm} publishes one metric of each kind. We report a case in which
a relation-scored statistic met every one of those grounds, along with six further checks
constructed against it, and the inference drawn from it was still wrong.

\paragraph{The control.} We constructed a second training corpus in which the relation the
metric targets has been destroyed while the quantities it might respond to instead --- amino-acid
composition and sequence length --- are preserved exactly. We pretrained the same architecture on
it under the same budget, fitted the same dictionary, and recomputed. The metric should not
systematically increase when the relation it depends on is absent from training.

\paragraph{Relation to prior work.} The manipulation itself is established.
\citet{sinha2021masked} pretrained masked language models on word-order-permuted text and found
that the resulting models scored surprisingly well on some parametric syntactic probes, which
they read as a deficiency in how representations are tested for syntactic information. We apply
the manipulation to an interpretability metric rather than to downstream capability. The nearest
analogue in the probing literature is the control task of \citet{hewitt2019designing}, which
randomises the labels to diagnose the probe's capacity; the control proposed here leaves the
labels alone and destroys the relation in the training distribution. The nearest existing check
in the SAE literature is the randomly initialised model: \citet{simon2025interplm} run a
weight-randomised baseline, \citet{heap2025automated} showed that aggregate interpretability
scores often fail to separate trained from random transformers, and \citet{korznikov2026sanity}
develop related checks. We run that check alongside ours, and \S\ref{sec:randominit} establishes
that the two diagnose different failures.

\paragraph{The specimen.} Throughout, we use a comparison between two training objectives,
masked and causal, on a controlled pair of pLMs: one backbone instantiated twice, sharing
initialisation seed, corpus and batch order, with the attention mask and the loss built from it
as the only difference. Existing cross-model comparisons of SAE features in pLMs vary objective,
architecture, data and scale together, so feature-level differences cannot be attributed to any
one factor \citep{gao2026structural}; the closest matched comparison equalises FLOPs or
validation loss across separately trained models \citep[Appendix~D]{cheng2024compute} and
evaluates downstream performance, whereas we match the training trajectory and use the pair to
test a metric. Here the objective is the sole candidate explanation, which makes the metric's
failure attributable rather than ambiguous. Section~\ref{sec:discussion} states what we do and do
not conclude from the comparison itself.

\paragraph{Contributions.} We introduce a training-distribution control for metrics that score a
relation. We build a metric of our own, show that it survives every standard check we can
construct, and then invalidate it from five directions: the corpus control, a random-initialisation
baseline, a no-model layer, direct probes, and a definitional sweep. No single one would have
caught every published metric we tested. We also show that the sequence-separation cutoff
determines the sign of the masked-versus-causal comparison, and that no value of it we tested
separates real from permuted training data.

\section{Setup}

\subsection{The pair}
\label{sec:pair}

We pretrained a 42.0M-parameter transformer in two arms: 30 pre-LN blocks, $d_{\mathrm{model}} =
320$, 5 heads, SwiGLU, RoPE, and tied input and output embeddings, with the block shape anchored
to ESM-C (EvolutionaryScale). One backbone was instantiated twice, and the only difference
between arms was the attention mask and the loss built from it; initialisation seed, corpus and
batch order were shared. The masked arm used 15\% masking with the standard 80/10/10 split; the
causal arm predicted the next residue. Both arms were pretrained on three million sequences from
the \texttt{ConvergeBio/uniref50} HuggingFace
stream (shuffled with seed 42 and a 50{,}000-sequence buffer, first three million sequences
retained)\opentag{f:uniref} at a maximum sequence length of 512 and a batch
size of 32. Three seeds per arm gave six models, trained for 40{,}283 steps and approximately
660M tokens, or 15.7 tokens per parameter; a scaling arm ran to 500 tokens per parameter, a
32-fold increase in budget, \pend{at one seed per arm}{t:scaleseeds}.

The pair does not hold the number of supervised positions fixed: the masked arm receives a target
at 15\% of positions and the causal arm at every position. Equalising supervised positions instead
of tokens leaves the effect in place (\S\ref{sec:confounds}).

Dictionary fitting and evaluation used a separate set of 1{,}500 SCOPe 2.08 domains filtered at
40\% sequence identity \citep{chandonia2022scope}, for which experimental structures are
available. The feature matrix holds one row per residue, using the first 510 residues of each
domain, for 293{,}760 rows; residue labels were measured on the untruncated 295{,}240 residues.
Nine depths were sampled, indexed by block and by position among the sampled depths; both indices
are used throughout (Table~\ref{tab:depths}). A cell is one (arm, depth, seed) combination. The
corpus control (\S\ref{sec:shuffled}) ran at blocks 11, 14 and 18, at three seeds per arm.

\begin{table}[t]
\caption{The nine sampled depths, by block index and by relative position in the stack.}
\label{tab:depths}
\begin{center}
\begin{tabular}{lccccccccc}
\toprule
Block & 0 & 4 & 7 & 11 & 14 & 18 & 22 & 26 & 29 \\
Depth & 0\% & 13\% & 25\% & 38\% & 50\% & 63\% & 75\% & 88\% & 100\% \\
\bottomrule
\end{tabular}
\end{center}
\end{table}

\subsection{The dictionary and the statistic}
\label{sec:statistic}

\Lstruct{} adapts the structural-clustering analysis of \citet{simon2025interplm}, which compares
a feature's mean activation over spatial neighbours against a within-protein permutation null. We
depart from their settings in two ways: a contact radius of 8\,\angstrom{} rather than
6\,\angstrom, following the CASP contact definition \citep{schaarschmidt2018assessment}, and a
minimum sequence separation of $|i - j| \ge 12$ rather than none, which excludes local backbone
contacts including those internal to $\alpha$-helices \citep{marks2011protein,morcos2011direct}
and retains the medium- and long-range classes on which contact prediction is conventionally
assessed. Contacts are measured between C$\alpha$ rather than C$\beta$ atoms.

The active set $A_f$ is the set of residues whose activation strictly exceeds the 90th percentile
of feature $f$'s activation across the evaluation set, and the structural neighbourhood is
%
\begin{equation}
\mathcal{N}(i) \;=\; \bigl\{\, j \;:\; \lVert \mathrm{C}\alpha_i - \mathrm{C}\alpha_j \rVert <
8\,\angstrom \;\; \text{and} \;\; |i - j| \ge 12 \,\bigr\},
\label{eq:neighbourhood}
\end{equation}
%
so, with $\bar{a}_f(i)$ the mean activation of $f$ over $\mathcal{N}(i)$ and
$\bar{a}^{\pi}_f$ the same quantity under a within-protein permutation $\pi$ of $f$'s
activations,
%
\begin{equation}
\Lstruct(f) \;=\;
\frac{\Bigl\{\mathbb{E}_{i \in A_f}\bigl[\bar{a}_f(i)\bigr] - \mathbb{E}_{i}\bigl[\bar{a}_f(i)\bigr]\Bigr\}
      \;-\; \mathbb{E}_{\pi}\Bigl[\mathbb{E}_{i \in A_f}\bigl[\bar{a}^{\pi}_f(i)\bigr]
            - \mathbb{E}_{i}\bigl[\bar{a}^{\pi}_f(i)\bigr]\Bigr]}
     {\mathrm{SD}\bigl(a_f\bigr)},
\label{eq:lstruct}
\end{equation}
%
with the null taken over five such permutations, and $\mathrm{SD}(a_f)$ the standard deviation
of $f$'s raw activations across the evaluation set. Both the observed and the null term are
centred by the mean of the neighbour-averaged activation over all residues; the two centring
terms do not cancel, because the permutation is within-protein while the neighbour average is
degree-weighted. Two implementation details matter for reading absolute values. With a TopK
dictionary the 90th percentile of a feature's activations is frequently exactly zero, so $A_f$
degenerates to the non-zero support rather than to a tenth of the residues; and any feature with
$|A_f| < 5$ is assigned $\Lstruct = 0$, so means taken over all 2{,}560 features include exact
zeros for dead and near-dead features. The statistic assumes that structural relations appear as co-activation of the
same sparse feature; \pend{information distributed across features, or encoded across pairs of
representations, would not be detected}{t:pairwise}. Section~\ref{sec:cutoff} varies the separation cutoff;
\S\ref{sec:dstruct} evaluates Simon and Zou's statistic at their own settings. The statistic is
not a Cohen's $d$: the numerator is in neighbour-averaged units and the denominator is the
standard deviation of raw activations.

\pend{Five permutations are few enough that the null estimate affects absolute levels. Raising it
to 25 permutations across 27 cells shifts mean \Lstruct{} by $+0.0073$, never by more than
$+0.0187$, and changes no cell's sign. Absolute levels below should therefore be read as
approximate; every between-condition comparison is unaffected.}{t:perm100}

Two symbols are used for between-arm comparisons, on different scales. $\Delta$ is the raw
masked-minus-causal difference in mean \Lstruct{} over all 2{,}560 features; the fold-disjoint
refit of \S\ref{sec:leak} gives $\Delta = +0.0141$. $d$ is the standardised between-arm effect
size over the three seed means, giving $d = +0.2568$ for the same comparison. The two differ by
roughly a factor of eighteen, and each figure below is labelled with the symbol it belongs to.

\subsection{The shuffled corpus}
\label{sec:shuffled}

We permuted residues within each sequence of the pretraining corpus at construction, independently
per sequence across the same three million sequences, from a single random stream. Length and
per-sequence amino-acid composition are preserved exactly, since a permutation reorders the same
multiset, while all order is destroyed, and with it the correspondence between sequence position
and spatial position. Both arms trained on the same shuffled corpus under the same budget;
dictionary fitting, evaluation set and the \Lstruct{} implementation were identical to the native
runs, and structural evaluation still uses real experimental structures.

Relative to their natively trained counterparts, causal held-out perplexity rises 10.6\% and
masked pseudo-perplexity 11.1\%. The near-equal degradation rules out one arm having learned the
shuffled corpus substantially better than the other.

We tested one destruction procedure, and a second permutation seed would not constitute a
different one: composition and length are fixed by construction, so any seed gives a statistically
near-identical corpus \pend{(\S\ref{sec:discussion})}{t:blockshuffle}.

\subsection{A leak in the evaluation split}
\label{sec:leak}

The dictionary was fitted on 1{,}350 domains and \Lstruct{} was scored on all 1{,}500, so every
fitting domain also appeared in the scored set. The original 150-domain holdout did not fix this:
118 of them (78.7\%) share a SCOPe fold with a fitting domain.

Scoring on those 150 domains alone reduces the effect by 30\% but leaves it significant ($d =
+0.1797$, 95\% CI $[0.1105, 0.2489]$). A fold-disjoint replacement, in which none of the 151
held-out domains shares a fold with a fitting domain, reproduces the arm separation to within
0.4\% across all 18 refitted cells ($\Delta = +0.01410$ against $+0.01415$). No structural
classification guarantees unrelatedness --- all extant proteins share evolutionary history --- but
fold-disjointness removes the level at which the leak operates.

The leak inflates the effect size but does not produce it. Figures below are computed on the
original split unless stated; Appendix~M gives the fold-disjoint values and the construction.

\section{Confounds that do not account for the effect}
\label{sec:confounds}

A bidirectional model attends to residues on both sides of a contact; a causal model carries
structural information forward from left to right. The natural prediction is that bidirectional
training produces representations in which structural co-activation is stronger, and \Lstruct{}
appears to confirm it.

\Lstruct{} is higher in masked than in causal models: $d = +0.2568$, seed means $0.2413 / 0.2455 /
0.2837$, 95\% CI over the three seed means $[0.1987, 0.3149]$. A two-sided $t$-test on the three
seed means (SD $0.02336$, SE $0.01349$, $t = 19.04$, $\mathrm{df} = 2$) gives $p = 0.0027$, and
seven of nine depths survive Benjamini--Hochberg correction. Under a prediction-matched protocol
that equalises the number of supervised positions instead of tokens, all 18 seed-by-depth cells
are significant.

We then constructed six checks, each testing whether the difference could arise without either
model having learned more structure. Depths are given by block index throughout.

\begin{itemize}[leftmargin=1.2em]

\item \textbf{Dictionary seed.} Is the effect dictionary-fitting noise? Across six arm-by-depth
cells ($2$ arms $\times$ 3 depths, one model seed) at three dictionary seeds, the median
seed-to-seed SD of mean \Lstruct{} is $0.00034$, against $0.00093$ across three model seeds
computed the same way on the same cells --- 37\% of it, and far below $\Delta = 0.0141$. Three
seeds make the SD estimate itself noisy.

\item \textbf{A frozen shared dictionary.} Do the two arms learn different dictionaries? Forcing
both arms onto one dictionary under a frozen decoder, a frozen encoder and a soft-frozen variant
\citep{korznikov2026sanity} leaves $\Delta$ positive in all nine conditions, across three depths
and three variants. It grows at the two shallower depths and shrinks at the deepest: block 7
$+0.0146 \rightarrow +0.0217$, block 14 $+0.0167 \rightarrow +0.0241$, block 22 $+0.0072
\rightarrow +0.0063$. A shared crosscoder \citep{minder2025crosscoders} gives the same direction
and decays with depth, at $+0.0209$, $+0.0160$ and $+0.0009$ for blocks 7, 14 and 22.
\pend{This condition runs at blocks 7, 14 and 22 rather than the 11, 14 and 18 of the corpus
control, sharing only block 14.}{t:frozendepths} The frozen conditions were run without a
bootstrap and carry no confidence intervals. Training details for all three variants and the
crosscoder are in Appendix~G.

\item \textbf{No dictionary at all.} Does the SAE create the effect? Recomputing \Lstruct{} with
no autoencoder anywhere, taking each of the 320 residual-stream dimensions as a feature in place
of the 2{,}560 SAE latents, leaves $\Delta$ at $+0.0084$, $+0.0056$ and $+0.0022$ at blocks 11, 14
and 18 --- positive at every depth, and in the same depth ordering as the SAE version.

\item \textbf{Reconstruction quality.} The masked arm's SAE reconstructs better than the causal
arm's (higher explained variance), which could inflate its \Lstruct. Across all 27 seed-by-depth
cells (9 depths $\times$ 3 seeds), the between-arm explained-variance gap correlates with the
effect at $\rho = +0.475$, but the correlation is carried entirely by the two shallowest depths
(blocks 0 and 4). At those depths the effect itself is near zero --- the two arms differ by
$0.00102$ at block 0 against \chknote{$0.0161$ or $0.01719$? reconcile with
Table~\ref{tab:untrained}}{c:recon} at block 11 --- and both fail Benjamini--Hochberg correction.
Over the remaining seven depths, $\rho = -0.043$ ($p = 0.85$). The confound is concentrated where
there is no effect to confound.

\item \textbf{Amino-acid composition.} Does the metric track which residues a feature fires on
rather than where they sit? For each feature we computed a background-corrected selectivity score
measuring how concentrated a feature's activations are across the 20 residue types
(Appendix~H). \pend{We predicted it would correlate positively with \Lstruct, since residues
sharing identity or chemical properties tend to cluster spatially. It does not: Spearman $\rho$
over all 2{,}560 features is $-0.044$ on the masked arm and $-0.122$ on the causal.}{t:top1share}

A control statistic confirms this is specific to the structural neighbourhood. \Lseq{} is
\Lstruct{} with the spatial neighbourhood replaced by the four nearest sequence neighbours $\{i-2,
i-1, i+1, i+2\}$, following InterPLM's sequential-clustering definition
\citep{simon2025interplm}: same formula, same null, same evaluation set, same features.
Selectivity correlates with \Lseq{} at $-0.676$ and $-0.412$. Composition drives the sequential
statistic strongly and the structural one weakly, and since the two differ only in the
neighbourhood, a confound acting through the dictionary or the evaluation set could not separate
them. Section~\ref{sec:indicator} establishes the complementary fact: the level a single
composition feature can reach is high.

\item \textbf{Bidirectional access.} The masked model attends to residues on both sides of a
contact; the causal model attends only to earlier sequence positions. This asymmetry could explain
the masked advantage if \Lstruct{} were driven by contacts whose partner lies later in sequence,
which the causal arm cannot see. That predicts the masked advantage should be larger for
later-in-sequence contacts than for earlier ones. We split each feature's structural neighbourhood
by contact direction and took the between-arm interaction, coded so that a positive value favours
later-in-sequence contacts (Appendix~K). It is $-0.1422$, negative at all nine depth-by-seed cells
(3 depths $\times$ 3 seeds), $p = 0.004$. The sign is the opposite of the prediction, so the
account is contradicted rather than merely unsupported.

\end{itemize}

With the caveats above, the effect survives all six checks. The masked advantage holds across
three model seeds, survives multiple-comparison correction at seven of nine depths, and orders the
two arms consistently. Section~\ref{sec:invalid} is what happens when the result is asked one
further question.

\section{Evidence that the inference is invalid}
\label{sec:invalid}

Five lines of evidence invalidate the inference behind the masked advantage; the first four are
tested independently, and the fifth concerns a parameter choice. A positive control on sequential
locality (\S\ref{sec:order}) confirms the machinery works. Section~\ref{sec:dstruct} shows that a
published statistic fails the same control, and \S\ref{sec:threechecks} compares three checks
across six metrics.

\subsection{It rises when order is removed from training}
\label{sec:order}

\Lstruct{} is higher after order-destroyed training in every combination of arm, depth and seed:
18 of 18 (2 arms $\times$ 3 depths $\times$ 3 seeds). At the model level that is six of six
independent pairs, and an exact one-sided sign test gives $p = 0.016$. One cell in the scaling arm
(masked, block 14) reverses; it lies outside this grid.

Averaged over three seeds at blocks 11, 14 and 18, the causal arm moves from $0.00240$ native to
$0.05201$ order-destroyed, and the masked arm from $0.01655$ to $0.02730$. The causal base is near
zero, so ratios are large and variable; the direction is what matters, and all 18 differences are
positive.

Increasing the budget 32-fold (from 15.7 to 500 tokens per parameter), the order-destroyed
\Lstruct{} moves by $-4$\% to $+23$\%, while native \Lstruct{} shifts by up to 116\%. The
order-destroyed level is largely insensitive to training budget --- a training-invariant floor.
Appendix~C gives all scaling-arm values.

\Lseq, the same statistic applied to sequential neighbours (offsets $\pm 1$ and $\pm 2$), favours
the causal arm at every depth --- the mechanistically expected result, since a model that reads
left to right should carry more sequence-local structure. \pend{Under shuffled-sequence
retraining, \Lseq{} in the masked arm collapses to 13--26\% of its native value at the five
deepest blocks}{c:lseq}: destroying sequence order destroyed sequential locality, and the measure
registered it. \Lstruct{} rose in every cell under the same intervention. The two statistics use
identical machinery on identical models; one responds correctly and the other responds backwards.
The positive control is partial --- \pend{the causal arm's \Lseq{} inflates 1.7--1.9$\times$
rather than falling}{c:lseq}, and we have no account of that asymmetry. It establishes that
co-activation statistics of this kind can detect a real signal and that the failure in \Lstruct{}
is specific, not generic. Full values are in Appendix~D.

No setting of the separation cutoff we tested separates native from order-destroyed \Lstruct;
\S\ref{sec:cutoff} gives the full sweep.

\subsection{Training does not lift the metric above its random-initialisation level}
\label{sec:randominit}

If \Lstruct{} measures learned structure, training should raise it above the level of an untrained
network. \citet{simon2025interplm} run a weight-randomised baseline on their own models; we run a
random-initialisation check on the controlled pair (Table~\ref{tab:untrained}).

\begin{table}[t]
\caption{Mean \Lstruct{} for trained and untrained networks at the three control depths. Trained
values are the mean over three model seeds with the seed-to-seed SD; the untrained baseline is the
mean over three initialisations (seeds 42, 43, 44).}
\label{tab:untrained}
\begin{center}
\begin{tabular}{llccc}
\toprule
Arm & Block & Trained ($n = 3$) & Untrained ($n = 3$) & Trained $\div$ untrained \\
\midrule
Masked & 11 & $0.01881 \pm 0.00122$ & $0.01821 \pm 0.00078$ & 103.3\% \\
Masked & 14 & $0.01667 \pm 0.00107$ & $0.01778 \pm 0.00014$ & 93.7\% \\
Masked & 18 & $0.01417 \pm 0.00070$ & $0.01765 \pm 0.00065$ & 80.3\% \\
\midrule
Causal & 11 & $0.00162 \pm 0.00034$ & $0.01790 \pm 0.00005$ & 9.0\% \\
Causal & 14 & $0.00144 \pm 0.00085$ & $0.01846 \pm 0.00052$ & 7.8\% \\
Causal & 18 & $0.00413 \pm 0.00101$ & $0.01901 \pm 0.00009$ & 21.7\% \\
\bottomrule
\end{tabular}
\end{center}
\end{table}

An untrained network scores about $+0.018$, flat across both architectures and all three blocks.
The trained masked arm never clears that floor --- it is within one seed-to-seed SD of it at block
11 and falls to 80\% of it by block 18. The trained causal arm sits an order of magnitude beneath
it throughout. The gap between arms is a training effect; the level is not.

This reverses the reading from \S\ref{sec:order}. It is not that masked training organises
structure; it is that causal training suppresses a property the architecture already has at
initialisation. The corpus control cannot distinguish these --- it shows the metric rises on
order-destroyed data, but not whether the rise is above or below the random floor. The two checks
diagnose different failures, which is why both are needed.

\red{[Figure~1 --- to be produced. Mean \Lstruct{} at blocks 11, 14 and 18. Filled markers are
trained values averaged over three model seeds (error bars: seed-to-seed SD); open markers are
untrained values averaged over three initialisations. The grey band spans the untrained cells
($+0.01765$ to $+0.01901$).]}\opentag{fig:one}

The same baseline separates the three structural measures reported here. Concept-F1 clears it
cleanly at the per-concept level, with several concepts at $1.000$ trained against $0.000$
untrained. Structure probes clear it partially: untrained AUROC runs $0.65$ to $0.71$ against
trained $0.65$ to $0.90$, so the bottom of the trained range sits at the untrained floor while the
top is well clear of it. \Lstruct{} does not clear it at all.

\subsection{Direct probes favour the causal arm}
\label{sec:probes}

\Lstruct{} scores SAE features; the underlying question is whether the model encodes structure at
all. Linear probes for helix, strand and burial, fitted directly on the residue representations
with no SAE anywhere, answer that question independently. Values are masked minus causal, averaged
over three seeds; helix, strand and burial are test F1, contact is precision at $L/5$
(Table~\ref{tab:probes}).

\begin{table}[t]
\caption{SAE-free probe results at three depths. Values are masked minus causal, so a negative
number favours the causal arm.}
\label{tab:probes}
\begin{center}
\begin{tabular}{lcccc}
\toprule
Depth & Helix & Strand & Burial & Contact P@$L/5$ \\
\midrule
Block 7 (25\%)  & $-0.0234$ & $-0.0442$ & $-0.0087$ & $+0.0051$ \\
Block 11 (38\%) & $-0.0512$ & $-0.0815$ & $-0.0126$ & $+0.0017$ \\
Block 14 (50\%) & $-0.0388$ & $-0.0697$ & $-0.0031$ & $-0.0058$ \\
\bottomrule
\end{tabular}
\end{center}
\end{table}

All three residue-label probes favour the causal arm at every depth and in every seed: 27 of 27 (3
depths $\times$ 3 labels $\times$ 3 seeds). Contact precision at $L/5$ does not follow them --- it
favours the masked arm at blocks 7 and 11 and the causal arm at block 14 --- so the 27-cell count
covers the three residue-label probes only. \citet{cheng2024compute} find a masked advantage on
contact prediction at 3B parameters and 200B tokens; our budget is three orders of magnitude
smaller, and we read this column as underpowered rather than as disagreeing with them.

\pend{Blocks 11 and 14 are two of the three depths at which the corpus control was run, and are
where \Lstruct{} gives its largest masked advantage. The probes therefore disagree with the metric
over two of its three control depths.}{t:probe18}

The direct measure of structural content disagrees with the indirect one at the depths the claim
rests on. The disagreement is not a logical contradiction, since \Lstruct{} measures spatial
clustering of SAE features while the probes measure linear decodability. Replacing the linear
probe with an MLP made all nine depth-by-label comparisons stronger and none weaker, with the
per-seed count still 27 of 27. Probe architectures are in Appendix~E.

\subsection{A single amino-acid indicator outscores almost every learned feature}
\label{sec:indicator}

\citet{simon2025interplm} report that SAEs trained on weight-randomised pLMs still yield hundreds
of features tied to individual amino acid types, while capturing no biological concepts at all. We
ask what such a feature scores on a structural metric, using a no-model baseline: a synthetic layer
of 29 indicator functions over residue identity and its groupings, comprising 20 amino acids, 8
standard biochemical groupings such as hydrophobic and aromatic (Appendix~I), and one all-ones
column. Each indicator is substituted for a feature's activation vector in
Equation~\eqref{eq:lstruct} and scored by the same code, with the same null and the same
evaluation set.

Scored this way, \texttt{aa:C}, the indicator for cysteine, reaches $+0.4221$, roughly four
times the best-scoring learned feature on the causal arm. It outscores every learned feature in
five of six cells (2 arms $\times$ 3 depths), and in the sixth exactly one feature beats it.

Cysteine is the expected winner, and the runners-up say why. Disulfide bonds place cysteine pairs
in spatial contact at arbitrary sequence separation; the next seven indicators are the hydrophobic
and aromatic residues and their class groupings, which cluster in three dimensions through core
packing. Both are chemical facts about proteins that hold without any model. Because an indicator
does not depend on the model, these values are identical in every cell: on the causal arm, five of
the 29 constants outscore all 2{,}560 learned features.

This establishes a level, not a mechanism. A quantity computable from residue identity alone
reaches most of the range the metric reports. Section~\ref{sec:confounds}'s composition check
establishes the complementary fact: the level a single composition feature can reach is high. So
the metric's scale is reachable by composition, even though its variation is not driven by it. A
metric with that property cannot be read as a measure of learned structure on magnitude alone.

\subsection{The sign depends on a definitional choice}
\label{sec:cutoff}

The masked advantage in \Lstruct{} exists only above a minimum sequence separation of $|i - j| \ge
5$. Sweeping the cutoff at block 11, the between-arm effect size moves from $d = -2.2096$ at gap 1
to $d = +0.3425$ at gap 12, crossing zero at gap 5 and plateauing past gap 6; deeper blocks cross
earlier, at gaps 4 and 3. Below gap 5, the causal arm wins. The cutoff therefore determines which
arm the metric favours, not merely by how much. Gap 12 sits in a wide flat region with the sweep
maximum at gap 24 or beyond, so it was not tuned to the result. But the result requires a
choice, and a different choice reverses it. The sweep is computed at a 6\,\angstrom{} contact
radius (Appendix~F gives all 48 rows); \S\ref{sec:order} uses the same sweep for a different
contrast.

The reversal has an arithmetic explanation. At 6\,\angstrom{} with no separation floor, consecutive
residues ($\approx 3.8$\,\angstrom{} apart) qualify as structural neighbours: 99.5\% of adjacent $(i, i+1)$
pairs fall inside the cutoff, and offsets $|i-j| \le 3$ account for 72\% of all edges. The
structural set substantially contains its own sequential baseline. \pend{At those settings
\Lstruct{} reproduces \Lseq{} to within 2.5\% at every depth}{c:agreement} --- the two are
measuring the same thing. The sign reversal is therefore the sequential result re-reported through
a structural definition that has absorbed the sequential baseline.

\red{[Figure~2 --- to be produced. Line plot: separation cutoff on the $x$-axis, $d$ on the
$y$-axis, one line per block, horizontal rule at zero.]}\opentag{fig:two}

\subsection{A published structural statistic under the corpus control}
\label{sec:dstruct}

InterPLM \citep{simon2025interplm} publish a structural-clustering statistic \dstruct{} addressing
the same operational question as ours: for each feature, they identify the highest-activation
residue per protein, compute the mean activation of residues within 6\,\angstrom{} of it, and
compare against a null of five within-protein permutations, sampling up to 100 proteins per
feature. The released repository does not implement \dstruct, so we reimplemented it from the
published Methods. One decision left open in that description changes the result; we report both
settings.

The procedure gates features at $0.6$ of a maximum activation. The Methods leave the reference
maximum open; the released normalisation code computes a dataset-wide per-feature maximum, which we
treat as the primary reading. Under a per-protein maximum the gate admits a median of 1{,}485 of
1{,}500 proteins per feature; under the dataset-wide maximum, about 98. The fifteen-fold difference
propagates into every number the statistic produces (Appendix~P).

\begin{table}[t]
\caption{\dstruct{} under the per-protein gate, aggregated over each arm. ``Significant features''
is the mean count of features called significantly structural.}
\label{tab:dstruct}
\begin{center}
\begin{tabular}{lcccc}
\toprule
& \multicolumn{2}{c}{Mean \dstruct} & \multicolumn{2}{c}{Significant features} \\
\cmidrule(lr){2-3}\cmidrule(lr){4-5}
Arm & Real & Shuffled & Real & Shuffled \\
\midrule
Causal & $0.4484$ & $0.4672$  & $985.7$ & $1349.3$ \\
Masked & $0.1376$ & $-0.0431$ & $323.3$ & $25.0$ \\
\bottomrule
\end{tabular}
\end{center}
\end{table}

On the causal arm, models trained on permuted residues score higher on both the mean and the count
of features called significantly structural, 1{,}349 against 986 (Table~\ref{tab:dstruct}). The
masked arm runs in the opposite direction. Under the dataset-wide gate the causal arm shows the
same direction block by block (\S\ref{sec:threechecks}, and Appendix~P, Table~2), at a higher
absolute level: $+0.5326$, $+0.5800$ and $+0.5800$ native against $+0.6086$, $+0.7475$ and
$+0.8713$ order-destroyed. The two gates disagree on the level and agree on the
failure, so the result does not turn on the ambiguity. Absolute values are not comparable across
gates and are not compared here. \pend{These values are a single seed.}{t:dstructseeds}

\dstruct{} is computed at 6\,\angstrom{} with no separation floor, which is the setting at which
\S\ref{sec:cutoff}'s cutoff sweep shows \Lstruct{} has no training-invariant floor. The two
statistics therefore fail the control at different settings, and the family resemblance is in the
construction rather than in the failure.

InterPLM use \dstruct{} to filter which features merit interpretation, not to compare models; they
report it once, on one model at one layer. Our result does not contradict any published comparison,
since none exists. It bounds what the statistic can be read as measuring.

Their second metric, concept-F1, scores how well individual SAE features align with known
biological labels. We evaluated it against 80 labels: 74 from SCOPe structural taxonomy (fold,
superfamily, family and class) and 6 residue-level annotations (secondary structure and solvent
accessibility)\opentag{t:conceptswissprot}. It passes our corpus control --- order-destroyed scores fall to $0.27$--$0.42\times$
of native on both arms. However, native values are close to zero throughout ($0.009$--$0.062$), so
the pass has limited weight. \pend{The aggregate favours the causal arm at four depths
(50--88\%), contradicting the probes of \S\ref{sec:probes}. The contradiction is an artefact of
aggregation: the 74 taxonomy labels favour the causal arm and dominate the aggregate, while the 6
residue-level labels favour the masked arm at the same depths, consistent with the probes.
Taxonomy labels are protein-level properties --- detecting ``fold c.1'' requires a feature that
fires somewhere in the domain --- and no single label's best feature exceeds a mean test-F1 of
$0.058$.}{c:conceptdecomp} Appendix~J gives the full decomposition.

Section~\ref{sec:threechecks} extends this, running all three checks across \dstruct{} and four
other published metrics alongside \Lstruct.

\subsection{The three checks catch different metrics}
\label{sec:threechecks}

The corpus control, the no-model layer and the standard random-init check are not redundant with
each other. Table~\ref{tab:threechecks} runs all three over six metrics on identical inputs.

\Lstruct{} and \dstruct{} are defined above. The remaining four are: concept-F1, the fraction of
biological concepts matched to at least one SAE feature \citep{simon2025interplm}, run using their
released code; single-latent F1, a per-feature classification score against structural and
functional annotations \citep{adams2025mechanistic}; the co-activation--decoder-geometry
correlation $\rho$, which measures whether features that frequently co-activate also have similar
decoder weight vectors \citep{li2025geometry}; and $k$-sparse probing AUROC from SAEBench, which
trains a sparse linear probe on SAE features to predict token-level labels
\citep{karvonen2025saebench}. Implementation details are in Appendix~J.

\begin{table}[t]
\caption{Six metrics under three validity checks. A metric \emph{fails} if its value on the
modified condition is at or above its value on the native model in the cells we ran, and
\emph{passes} if the modified condition scores below it. \emph{No verdict} covers cells that were
not run, are undefined on that input, or are underpowered; the reason is given in each such cell.
Verdicts are read off the values themselves; where a cell reports a significance test, that test is
the source's own. Per-cell values are in Appendix~J.}
\label{tab:threechecks}
\begin{center}
\small
\begin{tabular}{>{\raggedright\arraybackslash}p{0.27\textwidth}>{\raggedright\arraybackslash}p{0.20\textwidth}>{\raggedright\arraybackslash}p{0.19\textwidth}>{\raggedright\arraybackslash}p{0.19\textwidth}}
\toprule
Metric (source) & Shuffled-corpus & No-model & Random-init \\
\midrule
\Lstruct{} (ours) & fails --- 18/18 cells & fails --- \texttt{aa:C} above all & fails --- at floor \\
\addlinespace
\dstruct{} \citep{simon2025interplm}, reimpl. & fails on causal, passes on masked & passes --- 0 survive correction & \pend{no verdict --- not run}{t:dstructuntrained} \\
\addlinespace
concept-F1 \citep{simon2025interplm} & passes both arms & \pend{fails --- 4 of 12 concepts}{c:concept12} & passes --- separates \\
\addlinespace
single-latent F1 \citep{adams2025mechanistic} & no verdict --- underpowered & no verdict --- underpowered & no verdict --- underpowered \\
\addlinespace
co-activation $\rho$ \citep{li2025geometry} & passes --- 18/18 cells & no verdict --- undefined & no verdict --- underpowered \\
\addlinespace
SAEBench \citep{karvonen2025saebench} & no verdict --- underpowered & passes --- lowest of all & passes --- 6/6 cells \\
\bottomrule
\end{tabular}
\end{center}
\end{table}

Three metrics have enough data for a verdict on both the corpus control and the no-model check.
\Lstruct{} fails both, and fails the random-init check as well --- the only metric caught by all
three. Concept-F1 passes the corpus control but fails the no-model check. \dstruct{} fails the
corpus control on the causal arm and passes on the masked, at a single seed. The remaining cells
are either unrun, undefined for that input, or underpowered, and are not read as failures.

The checks are not ordered by severity. The no-model layer scores 29 indicator functions that do
not depend on any model, so its values are identical whatever corpus a model was trained on; it can
catch a metric whose scale is reachable without a model, and can say nothing about what training
does. The corpus control varies only the training distribution and holds the scoring layer fixed,
so it can catch a metric that rises when the relation is destroyed, and says nothing about the
level. \Lstruct{} fails both, for different reasons.

The single-latent F1 row is an absence rather than a failure: the metric has no discriminative
power on this concept set for any of the five inputs, a finding about prevalence-dominated concept
metrics rather than about the models. Appendix~J decomposes the 83--91\% figure into its two modes.

\section{Discussion}
\label{sec:discussion}

\paragraph{What is identified.} The metric rises when the relation is removed from training, which
removes the inference that a higher \Lstruct{} means more learned structure.
Sections~\ref{sec:randominit} and \ref{sec:probes} narrow what it measures instead: the quantity
survives random initialisation and disagrees with direct decodability. The stronger claim, that it
measures nothing about structure, is not identified.

\paragraph{Cost.} The control requires one additional pretraining run per arm at the original
budget (approximately 0.8 GPU-hours on one RTX PRO 6000), and scales with the model. It should be
run once per metric on a small controlled pair rather than once per model. The no-model layer
requires no pretraining and caught concept-F1, which the corpus control cleared; the corpus control
caught \dstruct, which the no-model layer cleared. The three are complements, and the cheap ones do
not substitute for the expensive one.

\paragraph{Limitations: the control.}
\begin{itemize}[leftmargin=1.2em]
\item \pend{One destruction procedure (\S\ref{sec:shuffled}). Replicating the treatment requires a
structurally different destruction, such as a cross-sequence shuffle preserving only global
composition, or a block shuffle preserving local windows, rather than another permutation
seed.}{t:blockshuffle}
\item Training- and evaluation-distribution shift are separated. Order-destroyed models were
rescored on an order-permuted evaluation set (18 cells). The causal arm keeps essentially all of
its elevation ($-9\%$ at block 11, $+3\%$ at block 18); the masked arm loses two-thirds and falls
below the natively trained masked arm. The permuted evaluation set permutes the residue string
only: no permutation indices are recorded, and the scoring code reads coordinates by position, so
activations were scored in permuted-index order against the native contact graph. The run
therefore measures co-activation over a real contact graph whose residue identities have been
decoupled: a stronger manipulation than scoring on order-destroyed input, and one that bears on
\S\ref{sec:order} rather than on this limitation.
\item \pend{Coverage. The corpus control ran at blocks 11, 14 and 18, the three depths where the
masked advantage is largest. Early and late blocks are untested.}{t:depthgrid}
\end{itemize}

\paragraph{Limitations: the measurement.}
\begin{itemize}[leftmargin=1.2em]
\item \pend{The denominator is untested. \Lstruct{} divides by $\mathrm{SD}(a_f)$, and we have not
recomputed it with a fixed or rank-based denominator, though the components are computed separately
in the analysis code.}{t:denominator}
\item Activation scales differ between conditions. Order-destroyed dictionaries have activation
scales roughly an order of magnitude below their native counterparts. Those scale ratios do not
track the ratios in \Lstruct{} and move in opposite directions between arms, which is evidence that
the scale gap does not propagate through the denominator rather than a direct test of
scale-invariance.
\item Reconstruction quality tracks the metric. Varying only the SAE while holding the model fixed,
\Lstruct{} tracks reconstruction quality (Appendix~G). \pend{The reconstruction gap between native
and order-destroyed models is large enough that, depending on the fit, it could account for 20\% to
all of the rise in \S\ref{sec:order}}{c:dose} --- a range six configurations per arm cannot narrow.
\item Dictionary quality differs between conditions. Seven of the thirty order-destroyed cells
carry a degeneracy flag set by the SAE pipeline, six of them on the masked arm and one on the
causal. Order-destroyed dictionaries also reconstruct better than native ones (validation explained
variance $0.99$ against $0.89$--$0.98$ masked, $0.88$--$0.91$ against $0.67$--$0.77$ causal), so a
dictionary-quality account of the rise remains open. The larger failure is on the causal arm, whose
dictionaries carry one flag between them, and \S\ref{sec:indicator} shows the metric's level is
reachable with no dictionary at all.
\end{itemize}

\paragraph{Limitations: the comparison.}
\begin{itemize}[leftmargin=1.2em]
\item \pend{The fold-disjoint split. Only the arm levels of \S\ref{sec:leak} were recomputed on it.
The corpus control of \S\ref{sec:order} is computed on the original split, where the contamination
is common to both conditions being compared.}{t:folddisjoint}
\item \pend{\dstruct{} coverage. The control is a single seed, and no untrained baseline was run
for it.}{t:dstructuntrained}
\item Perplexity is provisional. The masked-arm pseudo-perplexities depend on a mask-token
identification corrected during analysis.
\end{itemize}

\paragraph{Limitations: scope.}
\begin{itemize}[leftmargin=1.2em]
\item One scale. The controlled pair is 42M parameters. \pend{Whether the training-invariant floor
persists at larger scales is untested.}{n:onescale}
\end{itemize}

\paragraph{Scope of the claims.} We claim that a metric can rise when the relation it names is
destroyed, that six checks spanning dictionary training, reconstruction quality, composition and
architecture did not detect this, that a random-init baseline and a data control diagnose different
things, and that one extra pretraining run makes the failure visible. A workshop version reported
the masked advantage as a positive result \citep{gao2026structural}; the controls here supersede
that claim. Nothing here claims that masked pretraining is better or worse than causal at anything
biological, nor that any metric here is worthless outside the conditions we ran.

\subsubsection*{Reproducibility statement}

The controlled pair is 30 pre-LN blocks, $d_{\mathrm{model}} = 320$, 5 heads, SwiGLU and RoPE,
trained with AdamW ($\beta = 0.9, 0.95$), decoupled weight decay $0.1$ with zero decay on
normalisation weights, gradients clipped to global norm $1.0$, learning rate $6 \times 10^{-4}$
with 500 warm-up steps and cosine decay to one tenth of peak, bfloat16 autocast, no gradient
accumulation, batch size 32, maximum sequence length 512. Within each seed the two arms share
initial parameter values, training corpus and batch order; only the attention mask and the loss
built from it differ. Dictionaries are TopK with $k = 256$ and expansion 8, trained with AdamW at
$5 \times 10^{-5}$, batch size 4096 residues, 60 epochs, final-epoch weights, no checkpoint
selection, AuxK over 64 auxiliary latents at coefficient $1/32$. Appendices~A and B give the full
configuration; Appendix~P gives reproduction notes for both reimplemented statistics.

\bibliography{references}
\bibliographystyle{iclr2027_conference}

\appendix

\ifdefined\CLEANCOPY\else
\section{Outstanding items}
\label{app:working}

\red{\emph{Not part of the paper. Red text in the main text points here.}}

\subsection{To verify}

\begin{enumerate}[label=C\arabic*., ref=C\arabic*, leftmargin=2.6em]

\item \label{c:recon} \textbf{Reconstruction bullet, \S\ref{sec:confounds}.} ``$0.0161$ at block
11'' against Table~\ref{tab:untrained}, which gives $0.01881 - 0.00162 = 0.01719$. Reconcile.

\item \label{c:lseq} \textbf{\Lseq{} values, \S\ref{sec:order}.} ``Collapses to 13--26\%'' and
``inflates 1.7--1.9$\times$'' are from v1, seed 42. Recompute on three seeds.

\item \label{c:agreement} \textbf{$\Lstruct \approx \Lseq$ agreement, \S\ref{sec:cutoff}.} Stated
as ``within 2.5\% at every depth'', but the v1 pairs ($-2.266 / -2.210$, $-0.508 / -0.479$,
$-0.025 / -0.008$) give relative gaps of 2.5\%, 5.7\% and 68\%. Recompute or restate.

\item \label{c:conceptdecomp} \textbf{Concept-F1 decomposition, \S\ref{sec:dstruct}.} The 80
labels (74 SCOPe taxonomy, 6 residue-level) reproduce exactly from \texttt{cache/scope\_40.fa} and
\texttt{cache/residue\_features.csv} at \texttt{--min-domains 10}. The per-depth split and the
$0.039$--$0.058$ test-F1 range are still v1.

\item \label{c:concept12} \textbf{Table~\ref{tab:threechecks}, concept-F1 no-model cell.} ``4 of 12 concepts'' is
inconsistent with the 80 labels of \S\ref{sec:dstruct}; 12 is the family-level slice alone.
Identify which run produced it.

\item \label{c:dose} \textbf{Dose--response range, \S\ref{sec:discussion}.} ``20\% to all''
against v1's 4\%--129\%.

\end{enumerate}

\subsection{Runs queued}

\begin{enumerate}[label=T\arabic*., ref=T\arabic*, leftmargin=2.6em]

\item \label{t:top1share} \texttt{top1\_share} agreement. \texttt{ONLY=1 bash RUN\_TIER1.sh}.
Seconds. A disagreement drops the check count from six to five.

\item \label{t:depthgrid} Corpus control at the other six depths.
\texttt{ONLY=1 bash RUN\_DEPTH\_GRID.sh}. 75 min.

\item \label{t:denominator} Fixed and rank-based denominator.
\texttt{ONLY=2 bash RUN\_TIER1.sh}. 45 min. The per-feature SD is not stored, so this is a second
pass over the activations, not a re-read.

\item \label{t:probe18} Probes at block 18. \texttt{ONLY=2 bash RUN\_DEPTH\_GRID.sh}. Under an
hour.

\item \label{t:perm100} 100-permutation null. \texttt{NSHUF\_HI=100 ONLY=2 bash run\_checks.sh}.
Replaces the 5-versus-25 caveat in \S\ref{sec:statistic}.

\item \label{t:folddisjoint} Fold-disjoint refit of the corpus control.
\texttt{FOLDDISJ\_APPLY=1 ONLY=3 bash run\_checks.sh}. The shuffled half is already run.

\item \label{t:dstructseeds} Three seeds for \dstruct. \texttt{ONLY=3 bash RUN\_TIER1.sh}. 12 h,
CPU; 6 h with \texttt{GATES=global}.

\item \label{t:dstructuntrained} \dstruct{} untrained baseline.
\texttt{ONLY=4 bash RUN\_TIER1.sh}. 6 h. The only ``not run'' cell in
Table~\ref{tab:threechecks}.

\item \label{t:pairwise} Pairwise contact probes over pairs of residue representations, SAE and raw, with
shuffled-label and separation-matched controls. \texttt{STAGE=11 bash RUN\_MUSTRUNS.sh}. 1 h, CPU.
Tests the assumption in \S\ref{sec:statistic} that structure shows up as co-activation of a single
feature: \Lstruct{} cannot see a relational encoding.

\item Feature labelling. \texttt{label\_features.py} is untracked on the GPU box, so it needs
committing rather than writing.

\end{enumerate}

\subsection{Not queued}

\begin{enumerate}[label=N\arabic*., ref=N\arabic*, leftmargin=2.6em]
\item \label{t:frozendepths} Frozen-dictionary depths at 11 / 14 / 18. Retrains the shared
dictionaries; days.
\item \label{t:conceptswissprot} Concept-F1 on Simon and Zou's Swiss-Prot set. About 68 h.
\item \label{t:blockshuffle} Block-shuffle destruction. Needs pretraining.
\item \label{t:scaleseeds} Extra seeds in the scaling arm, currently $n = 1$ (seed 42). About
200 GPU-h.
\item \label{n:onescale} One scale. The pair is 42M parameters.
\item v1's top-features contact reversal ($+0.0300$ from raw activations against $-0.0113$
through the top-1024 \Lstruct{} features). Possible appendix.
\end{enumerate}

\subsection{To look up}

\begin{enumerate}[label=L\arabic*., ref=L\arabic*, leftmargin=2.6em]
\item \label{f:uniref} UniRef50 release behind the \texttt{ConvergeBio/uniref50} snapshot.
\end{enumerate}

\subsection{Figures}

\begin{enumerate}[label=FIG\arabic*., ref=FIG\arabic*, leftmargin=2.6em]
\item \label{fig:one} Mean \Lstruct{} at blocks 11, 14 and 18: trained (filled, $\pm$ seed-to-seed
SD) against untrained (open), grey band $+0.01765$ to $+0.01901$. (\S\ref{sec:randominit})
\item \label{fig:two} Cutoff sweep: separation cutoff on $x$, $d$ on $y$, one line per block,
rule at zero. (\S\ref{sec:cutoff})
\end{enumerate}

\fi

\end{document}
